\section{Finite alignment of the models}\label{sec:alignment}

We prove Principle~\ref{pr:alignment}.  The argument separates a finite
coloring construction from its realization by integer variables.  The finite
construction fixes all its choices before the threshold and the model
complexity are known.  The integer realization uses the marginal cube
symmetries from the preceding two Sections; its passage from existence to
positive probability loses only the number of those finite choices.

Throughout this Section, a target is a block $B=T\cup\{i\}$ with
$\mathcal E_B\ne\varnothing$.  A pair at pivot $i$ means
$e=(B,A)$, where
\[
 B=T\cup\{i\},\qquad A=P\cup\{j\},\qquad
 \varnothing\ne P<T<\{j\}<\{i\}.
\]
All lists below are finite.  Their sizes may depend on $n,r,s$, but never
on $w$, $\tau$, or the complexity of the models.

The integer meaning of one transformation is elementary.  At a center
$tp$, where $t$ is a tail product and $p$ is the pivot variable, an
integer displacement $q$ need not be divisible by $t$.  Since
$(t,M)=1$, choose an integer $\rho$ with $t\rho\equiv q\pmod M$ and put
$\Delta=(q-t\rho)/M$.  Then
\[
 t(p+m\rho)+mM\Delta=tp+mq\qquad(m\in\Z).
\]
The change $p\mapsto p+m\rho$ supplies the correct residue modulo $M$;
the remaining change $m\Delta$ is in the progression index of the
model on that residue class.  We will realize these two changes by
a transformation of the model state.  A transformation chosen for
one target can move the states for other targets as well.  This is
why the finite construction uses words in possibly noncommuting
transformations and protects the comparisons already obtained for
earlier targets.  A general word need not translate the integer
variable $p$.  The precise residue choices and the states on which
these transformations act are constructed after the finite plan.

\subsection{A finite plan for words and rational scales}

Put $p_\ell=2^{n-\ell}$ and $p_I=\sum_{\ell\in I}p_\ell$.
At an update targeting $B_0=T_0\cup\{i\}$, define the multiplier tuple
$\kappa_{i,T_0}(k)$, for $k\in\N$, by
\begin{equation}\label{eq:alignment-scale-update}
 \kappa_\ell(k)=
 \begin{cases}
 k^{p_\ell},&\ell<i,\\
 k^{-p_{T_0}},&\ell=i,\\
 1,&\ell>i.
 \end{cases}
 \qquad b'=b\kappa(k).
\end{equation}
Here multiplication of tuples is coordinatewise.  Thus $b'_{B_0}=b_{B_0}$.
For every pair $e=(U,A)$ at this pivot, with $U=C\cup\{i\}$,
\begin{equation}\label{eq:alignment-positive-exponents}
 \frac{b'_A}{b'_U}=\frac{b_A}{b_U}k^{\lambda_e},
 \qquad \lambda_e=p_A-p_C+p_{T_0}>0.
\end{equation}
Indeed, if $a=\min P$, then
$p_a>\sum_{\ell>a}p_\ell\ge p_C$, so already $p_P>p_C$.
This proves positivity for every pair at the pivot, including those with
$C\ne T_0$.
These compensating rescalings follow the mechanism in
\cite[proof of Proposition~4.1]{Alweiss}.  Here the transformations
for different targets may not commute, so the comparisons that must
survive an update will be indexed by words.

For each pivot $i$, let a group $G_i$ of nilpotence class at most $s$
act on a set $\mathcal X_i$.  Associate to every pair $e=(B,A)$ at
that pivot an element $L_e\in G_i$.  For every target $B$ at that
pivot and every positive rational $a$, let
$C_{B,a}:\mathcal X_i\to\{0,1\}^r$ be an arbitrary map, called its
palette.  The finite choices below must work for every such collection
of actions and maps.

A word template is a word in letters $(e,\gamma)$ and their inverses,
where $\gamma\in\Q$.  At scale $b$ it is interpreted as a group word by
\begin{equation}\label{eq:alignment-letter}
 (e,\gamma)[b]=L_e^{q_0\gamma b_A/b_U},\qquad e=(U,A).
\end{equation}
The integer $q_0$ is chosen after the finite plan; all exponents actually
used will then be integers.  Words act on the left, so $wJx$ means first
apply $J$, then $w$.  The empty word is allowed.  At a center
$x\in\mathcal X_i$, a requirement for a block $B$,
leading-multiplier list $\mathcal L$, and word list $\mathcal W$ is
\begin{equation}\label{eq:alignment-word-requirement}
 C_{B,a b_B}(x)=C_{B,a b_B}(w[b]x)
 \quad(a\in\mathcal L,\ w\in\mathcal W).
\end{equation}

Nilpotent polynomial recurrence was developed by Leibman
\cite{Leibman1994} and by Bergelson and Leibman through polynomial
Hales--Jewett theorems \cite{BergelsonLeibman1999,BergelsonLeibman2003}.
We use the finite-coloring consequence of nilpotent polynomial recurrence
in \cite[Corollary~3.7]{ZorinKranich}.  In the form needed
here, for finitely many ordered words
\[
 w_j(k)=a_{j,1}^{P_{j,1}(k)}\cdots a_{j,l_j}^{P_{j,l_j}(k)},
 \qquad P_{j,l}\in\Z[k],\quad P_{j,l}(0)=0,
\]
in a nilpotent group of bounded step, every finite coloring admits $k>0$
and $J$ for which $J,w_1(k)J,\ldots,w_q(k)J$ have the same color.
Only the subgroup generated by the finitely many letters is involved;
it is countable, as required by that Corollary.  The recurrence statement
permits these noncommuting ordered products.
For the handedness, apply its left-translate formulation to the inverse
words and to the coloring $g\mapsto C(g^{-1})$, and then invert the
resulting pattern.  To check the polynomial hypothesis of that
specialization, let $D\ge1$ bound the degrees of the exponents and
stretch the lower central series to the filtration
$H_0=G$, $H_j=\gamma_{\lceil j/D\rceil}(G)$ for $j\ge1$.
The map $\alpha\mapsto a^{P(|\alpha|)}$ is an IP polynomial for
this filtration, since repeated derivatives are ordinary finite
differences of its exponent.  Product and inverse closure, as in
\cite[Theorem~2.5]{ZorinKranich}, includes the required ordered
words.  Include the identity word in the recurrence family.
Here $k=|\alpha|$ for a nonempty finite index set $\alpha$, so the
parameter is positive.  Stretching the filtration verifies the
recurrence hypothesis; the acting group still has nilpotence
class at most $s$.

There is a useful uniform finite version of this statement.  Give every
letter type its own generator in the free nilpotent group of step $s$ on
these finitely many generators.  On the compact space of its colorings
with a fixed finite palette, each successful pair $(k,J)$ defines an open
cylinder set. The preceding recurrence statement shows that these cylinders cover the
space.  A finite subcover gives finitely many pairs $(k,J)$.  Pulling a
coloring back by the homomorphism that interprets the generators proves
the same assertion in every nilpotent group of step at most $s$.  The
finite alternatives depend on the word templates, polynomial exponents,
palette size, and $s$, and on no numerical values assigned to the
generators.

\begin{lemma}[Finite word plan]\label{lem:finite-word-plan}
For given $n,r,s$, there are a positive integer $q_0$ and finite
lists of scale updates and word jumps with the following property,
uniformly over the actions and palette maps just specified and the
initial centers at all pivots.  Start the scale tuple at $(1,\ldots,1)$,
process the pivots increasingly, and fix an order of their targets.
At each pivot one can choose updates
\eqref{eq:alignment-scale-update} and jumps from these lists so that,
at the resulting scale $b$ and center, the terminal palette comparisons
hold simultaneously for all its targets.  These comparisons are
\eqref{eq:alignment-word-requirement} with leading multiplier $d_B$
and word
\begin{equation}\label{eq:alignment-terminal-word}
 \prod_{A\in D}\bigl((B,A),d_A/d_B\bigr)[b],
\end{equation}
in a fixed order, for every target $B$, every $D\subset\mathcal E_B$,
and every multiplier tuple $d$ obtainable from later prescribed updates.
When later pivots are processed, the center at this pivot is left fixed:
these terminal requirements already include all the later scale changes.

All these lists depend only on $n,r,s$.  There are finite lists
$\mathcal B\subset\Q_{>0}^n$ and $\mathcal A\subset\Q_{>0}$
containing, respectively, the final scale tuples and every palette index
consulted during the construction.  All intermediate scales belong to
finite lists as well.  The list $\mathcal A$ also contains every required
product of entries of a tuple in $\mathcal B$.  All powers used in
\eqref{eq:alignment-letter} are integral.  If each $L_e$ also translates
an auxiliary integer coordinate $v_e$ by $1$, all auxiliary indices
visited from $v=0$ lie in a prescribed finite slot list.
\end{lemma}

\begin{proof}
We give the backward recursion, including the requirements protecting
an earlier target.  At a fixed pivot let its ordered targets be
$B_1,\ldots,B_q$.  Suppose the required lists after update $j$ have
already been specified for $B_1,\ldots,B_j$; denote them by
$\mathcal L_{j,B}$ and $\mathcal W_{j,B}$.  At $j=q$ these are the
terminal lists in the statement.  Taking a full product of the leading
and word lists only strengthens the requirements.

Target $B_j$ at update $j$.  Its leading product does not change.
For a word $w$, let $I_k(w)$ denote the old-scale template obtained by
replacing each letter $(e,\gamma)$ by
$(e,\gamma k^{\lambda_e})$.  Equation~\eqref{eq:alignment-positive-exponents}
gives the exact identity
\begin{equation}\label{eq:alignment-word-inflation}
 w[b\kappa(k)]=I_k(w)[b].
\end{equation}
For $B_j$ and a fixed current center $x$, color the acting group by
$g\mapsto(C_{B_j,a b_{B_j}}(g x))_{a\in\mathcal L_{j,B_j}}$.
Its size is at most $2^{r|\mathcal L_{j,B_j}|}$.
Regard each rational old-scale letter needed in
$\mathcal W_{j,B_j}$ as a separate abstract generator.  The words
$I_k(w)$ are ordered products of these generators to powers
$\pm k^{\lambda_e}$, with all exponents vanishing at zero.
The finite version of recurrence therefore supplies a finite set
$\Omega_j$ of alternatives $(k,J)$ satisfying the required equalities
for $B_j$ at $J[b]x$ and scale $b\kappa(k)$.  This set is independent
of the incoming numerical scale.

For each previously processed $B=B_l$, $l<j$, define predecessor lists
by including, for every $(k,J)\in\Omega_j$,
\begin{align}
 \mathcal L_{j-1,B}&\supset
   \{a\kappa_B(k):a\in\mathcal L_{j,B}\},
       \label{eq:alignment-leading-budget}\\
 \mathcal W_{j-1,B}&\supset
   \{J\}\cup\{I_k(w)J:w\in\mathcal W_{j,B}\}.
       \label{eq:alignment-word-budget}
\end{align}
The unions over all alternatives are finite.  If these predecessor
requirements hold at $(x,b)$, then for each new leading index
$a b_B\kappa_B(k)$ they compare the same old center $x$ both with
$J[b]x$ and with $I_k(w)[b]J[b]x$.  Equality through the old center and
\eqref{eq:alignment-word-inflation} give
\[
 C_{B,a b'_B}(J[b]x)
   =C_{B,a b'_B}(w[b']J[b]x).
\]
These are exactly the requirements after the update.  Thus
\eqref{eq:alignment-leading-budget}--\eqref{eq:alignment-word-budget}
prove preservation, with no assumption that the jump fixes earlier
targets.  At $j=1$ there is no earlier target to protect.  Backward
induction defines the whole plan, and forward induction executes it.

Next plan the pivots in decreasing order.  Once the plans at pivots
larger than $i$ are fixed, take all coordinatewise products of their
possible multiplier tuples; call this finite list $\mathcal D_i$.
Use $\mathcal D_i$ as the list of $d$'s in the terminal requirements
at pivot $i$.  The preceding construction returns multiplier options
depending only on that list and the fixed data, so there is no
dependence on an unknown incoming scale.  The last pivot starts with
the singleton future list $\{(1,\ldots,1)\}$.

Starting all scales at $1$, enumerate all intermediate and terminal
paths through these finite alternatives.  This gives finite scale
lists and finitely many rational palette indices $a b_B$ at every
checkpoint.  Define $\mathcal B$ to contain the final scales and
$\mathcal A$ to contain all these indices and all needed subset
products of final scales.  Now choose $q_0$ divisible by the
denominators of every rational number $\gamma b_A/b_U$ used in any
evaluated word on any enumerated path.  This does not affect the
abstract recurrence choices: $q_0$ merely changes the homomorphism
interpreting each old-scale generator.

Finally enumerate all prefixes of all words, their cumulative jumps,
and the terminal tests, with the now integral exponents.  Their sums
in the auxiliary translation coordinates give a finite slot list,
which we enlarge to contain $0$.  Intermediate word evaluations use
only this list.  One may assign arbitrary palettes at other slots
when defining the coloring on the whole acting group.  Every step
of the construction used only $n,r,s$.
\end{proof}

\subsection{Integer residue shifts and polynomial arrays}

Fix a pivot $i$ having targets, and condition initially on the earlier
raw variables.  Write $p=t_i$, set $L=MW^w$, and, for every pair
$e=(U,A)$ at this pivot, write $U=C\cup\{i\}$ and define
\begin{equation}\label{eq:alignment-residue-data}
 q_e=\frac{h_A}{q_0h_U}t_A,\qquad
 t_Cr_e\equiv q_e\pmod L,\quad 0\le r_e<L,
 \qquad \Delta_e=\frac{q_e-t_Cr_e}{M}.
\end{equation}
These are integers for all sufficiently large $w$.  In fact
$h_A/h_U$ is a multiple of $W^w$, which is eventually divisible by
$q_0W$, and $t_C$ is a unit modulo the smooth modulus $L$.
Consequently
\begin{equation}\label{eq:alignment-residue-divisibility}
 W\mid q_e,\qquad W\mid r_e,\qquad W^w\mid\Delta_e.
\end{equation}
The last assertion follows from the full modulus $MW^w$ in the
congruence.  In particular every fixed positive integer eventually
divides $\Delta_e$.
For an integer slot vector $v$ put $s(v)=\sum_e v_er_e$.
The identity realizing one own-target letter is
\begin{equation}\label{eq:alignment-integer-identity}
 t_C(p+m r_e)+mM\Delta_e=t_Cp+m q_e\qquad(m\in\Z).
\end{equation}
It uses no divisibility of $q_e$ by $t_C$.

Let $D_i^-=\prod_{\ell<i}X_\ell^2$ and use microcells in the pivot
variable of length
\begin{equation}\label{eq:alignment-microcell}
 R=M\left\lfloor H_i^{1/2}/M\right\rfloor,
 \qquad N_{\rm loc}=R/M.
\end{equation}
Admissibility implies that $R$ dominates every fixed power of
$M+D_i^-$, that $D_i^-R/H_i\to0$, and that $R/X_i\to0$.
Since $W^w\le M$, all slot shifts and all bounded array evaluations
below have size at most a fixed power of $M+D_i^-$ and are $o(R)$
in the pivot variable.

After discarding partial endpoint cells, condition on a microcell and
on $r_0=p\bmod M$.  Its first point of that residue is $p_*$, and
\[
 p=p_*+M\ell,\qquad 0\le\ell<N_{\rm loc}.
\]
Because $r_0$ is a $W$-unit, every point of this progression is a
$W$-unit.  Its harmonic weights have ratio $1+O(R/X_i)$, so the
conditional law of $\ell$ differs in total variation by $o(1)$ from
the uniform law, uniformly in the earlier variables.

For $B=T\cup\{i\}$, put $t=t_T$.  For every prescribed slot and
every prescribed bounded integer array index $u=(u_e)_{e=(B,A)}$,
the numerical arguments we will use are
\begin{equation}\label{eq:alignment-array-argument}
 t(p+s(v))+M\sum_{e=(B,A)}u_e\Delta_e.
\end{equation}
Except on a set of conditionings of probability $o(1)$, all these
arguments, for the whole microcell, lie in one $H_i$-interval for
each $B$.  Here and below the constants may depend on the fixed
finite word lists.  To verify the boundary assertion uniformly, in
a dyadic pivot range $[Y,2Y)$ the preimages of $H_i$ boundaries have
$O(Yt/H_i+1)$ neighborhoods of length $O(R)$.  An interval of length
$a$ has $a\phi(W)/W+O(\phi(W))$ $W$-units.  Relative to the
$W$-units in $[Y,2Y)$, these neighborhoods therefore cost
\[
 O\left(\frac{t(R+W)}{H_i}+\frac{R+W}{Y}\right)=o(1).
\]
Harmonic weights change this bound by at most a constant on each
dyadic range.  Averaging the ranges and using $t\le D_i^-$ proves
the assertion.  Endpoint cells have still smaller cost.

On a good conditioning choose the nilsequence formula for
$S_{B,a,c}$ at the common interval and the residue
$r_v=t(p+s(v))\bmod M$.  For each $v,a,c$, write that formula as
$F(g^k x\Gamma)$, using a group representative $x$ and absorbing
the fixed origin of $k$ into $x$.  Retain the entire group-valued
polynomial array
\begin{equation}\label{eq:alignment-polynomial-array}
 u\longmapsto
 g^{(t(p+s(v))-r_v)/M+\sum_{e=(B,A)}u_e\Delta_e}x.
\end{equation}
Only finitely many evaluations must agree with the model; the array
itself is defined for every $u$.

We record why these arrays lie in fixed compact nilmanifolds.  If
$G$ has lower central Lie algebra series $\mathfrak g_j$, let
$\mathfrak a_q(G)$ consist of polynomials in $q$ real variables
whose constant coefficient is in $\mathfrak g$ and whose coefficient
of total degree $j>0$ is in $\mathfrak g_j$.
Pointwise brackets preserve this space: a bracket of degrees $j,k$
lies in $\mathfrak g_{j+k}$ when both are positive, and a bracket
with a constant coefficient lies in a still deeper allowed level.
It is a finite-dimensional rational nilpotent Lie algebra of step
at most $s$.  Its simply connected group $\mathcal P_q(G)$ consists
of the corresponding arrays, with pointwise multiplication.
\[
 \mathcal P_q(\Gamma)
   =\{a\in\mathcal P_q(G):a(z)\in\Gamma\text{ for every }z\in\Z^q\}
\]
is a lattice.  For discreteness, sufficiently many integer evaluations
determine every coefficient of the logarithm; a convergent sequence
of lattice-valued arrays is eventually constant at those evaluations
and hence constant.  For cocompactness, choose a rational basis
respecting the lower central series.  The change between logarithmic
coordinates and ordered Mal'cev coordinates is a rational polynomial
with zero constant term.  Clearing its finitely many denominators
shows that the exponential of a sufficiently divisible integral
coefficient vector is lattice-valued at every integer input.
In a rational Mal'cev basis for $\mathfrak a_q(G)$, fixed integer
multiples of the basis vectors therefore exponentiate into
$\mathcal P_q(\Gamma)$.  Successively reducing the ordered
coordinates modulo these elements, starting with the quotient by
the next ideal in the Mal'cev flag, places every coset in a bounded
coordinate box.  This proves cocompactness.

Integer input shifts $u\mapsto u+\mathbf e$ preserve the Lie algebra
and this lattice, and evaluation at a fixed integer input induces
a smooth map of the quotients.  Formula~\eqref{eq:alignment-polynomial-array}
belongs to this group: in its BCH logarithm a term of degree $j$
in $u$ contains at least $j$ occurrences of $\log g$, and hence
lies in $\mathfrak g_j$.

For each target take the product over all $v,a,c$ of these array
quotients, denoting it by $K_B/\Lambda_B$, and put
\[
 K/\Lambda=\prod_B K_B/\Lambda_B.
\]
The choices range over a fixed finite list because the original
model complexity is bounded.  Increasing $\ell$ by $1$ multiplies each
array on the left by the constant array $g^t$.  Thus the joint state
is a linear orbit on $K/\Lambda$, adapted to its lower central
filtration and of step at most $s$.  Denote its state at $\ell$ by
$Y(\ell)$.  For an own pair $e=(B,A)$ write $\mathcal S_e$ for the
lattice-preserving automorphism of the whole component $K_B$
shifting the corresponding input coordinate~by~$1$.

\subsection{Typical marginal cube invariance}

The state $Y(\ell)$ now stores every model reading needed at this pivot.
We must show that, for typical conditionings, its marginal cube laws
permit the array shifts $\mathcal S_e$.  Removing the rough tail factors
from their sampling steps will supply these symmetries.

\begin{lemma}[Conditional face invariance]\label{lem:conditional-face-invariance}
Fix the bounded-complexity model family and the finite lists above.
For a good conditioning on the earlier variables, the pivot microcell,
and its residue, form the joint orbit $Y(\ell)$ just constructed.
Fix a cube dimension $a\le s+1$, a positive integer $d$, constants
$0<c<C$, and a bound for the supremum and Lipschitz norms of tests.
Uniformly over boxes for $(\ell_0,\ldots,\ell_a)$ in
$[-CN_{\rm loc},CN_{\rm loc}]^{a+1}$ with sides at least
$cN_{\rm loc}$, all fixed residue vectors modulo $d$, all such tests
of one marginal cube, and every own-pair automorphism, the difference
between the marginal cube average and its image under that
automorphism on an upper codimension-one face tends to zero in
probability over the conditioning.

Consequently, from any sets of conditionings of probability bounded
away from zero one can select a sequence on which, simultaneously
for all these fixed data, the marginal face invariances hold.  On
a further subsequence to which Proposition~\ref{prop:local-cube-law}
applies, every projected local cube coset has the corresponding face
symmetry in dimensions through $s+1$.  The modulus in this last
assertion may be the period selected by that Proposition.
\end{lemma}

\begin{proof}
Fix a target $B=T\cup\{i\}$ and write $t=t_T$.  Formulae outside
the conditioned cell mean the extensions of the already selected
nilsequence formulae, so all bounded-multiple boxes in the statement
are defined.  We will compare them with averages whose step has no
factor $t$.

\paragraph{An exposure retaining polynomial dependence.}
We prove the assertion in probability by changing the order of exposure.
We must freeze the target's model formula while leaving its tail
variables available for progression comparison.  A fixed pivot microcell
alone does not do this as the tail varies.  Instead we condition on a
short bin for $tp$;
the eventual microcell origin will enter the base coordinate $x_0$,
not a coefficient of the polynomial family.
Expose the earlier variables outside $T$, the residues modulo
$L=MW^w$ of those inside $T$, dyadic bins $[S_\ell,2S_\ell)$ for
$\ell\in T$, the pivot residue $r_0\bmod M$, and a bin for $tp$ of
width $t_{\rm lo}R$, where $t_{\rm lo}=\prod_{\ell\in T}S_\ell$.
Do not expose the exact pivot microcell at this stage.
Then $t_{\rm lo}\le t\le2^{|T|}t_{\rm lo}$.  Every $r_e$ is
fixed: earlier residues determine $q_e\bmod L$ and $t_C\bmod L$.
For own pairs $e=(B,A)$ the full integer $q_e$ is fixed, because
$A\cap T=\varnothing$.  Hence all slot shifts $s(v)$ are fixed.
So are
\[
 r_*=tr_0\bmod M,\qquad r_v=t(r_0+s(v))\bmod M.
\]

We may discard exposed bins whose entire preimage in $p$, for every
$t$ in the indicated product of dyadics, is not inside
$[X_i+2R,X_i^2-2R]$.  Such a discard forces the actual $p$ into a
fixed multiplicative neighborhood of a cutoff endpoint, up to an
$O(R)$ enlargement.  Its harmonic probability is
$O_n(1/\log X_i)+o(1)$.  Also discard $tp$ bins within a sufficiently
large fixed multiple of $t_{\rm lo}R$ of an $H_i$ boundary.  The
boundary count preceding the Lemma gives probability $o(1)$ for
this discard.  The constant is chosen to cover the given
bounded-multiple boxes, all slots, and all required bounded array
indices.  On the remaining bins the interval and residue selecting
every formula for this target are fixed throughout the exposure.
Thus all its parameters $F,g,x$ are fixed as the variables in $T$
vary.

For later use, the conditional law of those variables is dominated
by a constant, depending only on $n$, times the product of uniform
laws on their dyadic bins and prescribed $L$ residues.  Before
conditioning on the $tp$ bin, their harmonic densities in each
dyadic differ from uniform by at most a factor $2$ per variable.
If the exposed bin is $[Q,Q+t_{\rm lo}R)$, its preimage at a given
$t$ has length $t_{\rm lo}R/t$, between $2^{-|T|}R$ and $R$.
It lies fully inside the pivot cutoff.  On the residue $r_0\bmod M$
the unnormalized pivot mass there is
\[
 (1+o(1))\frac1M\int_{Q/t}^{(Q+t_{\rm lo}R)/t}\frac{du}{u}
 = (1+o(1))\frac1M\log\left(1+\frac{t_{\rm lo}R}{Q}\right),
\]
uniformly in $t$.  The discretization error is relatively
$O(M/R)$, and the common normalization cancels.  In particular
conditioning on this bin changes the preceding product comparison
by a bounded factor.  This is the required domination; no exact
microcell has been fixed in deriving it.

\paragraph{Putting the tail product in the sampling step.}
Let $p_*$ be the first point of $r_0\bmod M$ in the actual
microcell eventually selected.  Put $Y_*=\lfloor Q/M\rfloor$ and
change the base and increment coordinates by
\begin{equation}\label{eq:alignment-cube-coordinates}
 x_0=\frac{t(p_*+M\ell_0)-r_*}{M}-Y_*,
 \qquad x_j=t\ell_j\quad(1\le j\le a).
\end{equation}
The unknown cell origin has entered $x_0$ rather than a polynomial
coefficient.  Since $p_*$ is within $R$ of the actual pivot, these
coordinates lie in boxes in a fixed multiple of
\[
 Z=\frac{t_{\rm lo}R}{M},
\]
with all sides comparable to $Z$.  The original residue restrictions
on $\ell_0,\ldots,\ell_a$ become one residue vector modulo $dt$ in $x$.
Changing an endpoint by less than $dt$ has vanishing relative
counting cost, so these are precisely the permitted normalized
box averages.

At vertex $\omega\in\{0,1\}^a$, set
$X=x_0+\sum_{j=1}^a\omega_jx_j$.  Its array exponent is
\begin{equation}\label{eq:alignment-joint-polynomial-exponent}
 X+Y_*+\frac{r_*+t s(v)-r_v}{M}
       +\sum_{e=(B,A)}u_e\frac{q_e-tr_e}{M}.
\end{equation}
Every parameter here except $X$ and $t$ was fixed by the exposure.
The displayed exponent is a polynomial of degree at most two,
with its only possible quadratic terms of the form $tu_e$.
BCH then shows that the logarithms of the arrays,
viewed in the fixed array group, are ordinary polynomials jointly
in $x$ and the variables $t_\ell$, $\ell\in T$, of bounded degree.
There is no bound required on their coefficients.  This joint
polynomiality is the hypothesis that permits Proposition~\ref{prop:rough-step}.

\paragraph{Removing each factor of the step.}
We compare the sampling in a class modulo $dt$ with the sampling
in the same class modulo $d$, removing the factors of $t$ in a
fixed order.  Consider removal of $t_\ell$, and condition on all
other factor values.  Let $w'$ be the product of the factors still
in the step after this removal.  The current step is $dw't_\ell$.
There is a representative $a_0$ for its base coordinate modulo
$dw'$ which can be chosen independently of $t_\ell$, after
allowing all the finitely many residues modulo $d$.
Indeed the base congruence modulo $w'$ is
\begin{equation}\label{eq:alignment-origin-congruence}
 M(x_0+Y_*)+r_*\equiv0\pmod {w'},
\end{equation}
and the increment congruences are $x_j\equiv0\pmod {w'}$.
Since $w'$ is a product of $W$-units,
$(M,w')=1$, and also $(d,w')=1$ for large $w$.
The Chinese remainder theorem therefore gives representatives
$a_j\in[0,dw')$ determined by these congruences and the chosen
$d$ residues, independently of $t_\ell$.

Substitute $x_j=a_j+dw'y_j$.  The current restriction becomes a
single class vector modulo $t_\ell$ in $y$, while removing it gives
unrestricted integer $y$.  No coprimality of $w'$ and $t_\ell$ is
used: both original classes are nested congruences, and division
of $x_j-a_j$ by the integer $dw'$ is exact.
The new common box scale is $Z/(dw')$.  It dominates every fixed
power of $S_\ell$, uniformly over the other factor values, by
the choice of $R$.  The substitution has coefficients independent
of $t_\ell$, so joint log-polynomiality in $(y,t_\ell)$ is retained.
The variable $t_\ell$ lies in a fixed residue modulo $L=MW^w$
in $[S_\ell,2S_\ell)$, with
\[
 S_\ell/L\longrightarrow\infty.
\]
All hypotheses of Proposition~\ref{prop:rough-step} hold.

That Proposition bounds the exceptional proportion uniformly in
the polynomial coefficients, all permitted boxes, class vectors,
and tests.  Its uniformity in endpoints is essential: the exact
microcell and hence the endpoints may depend on the pivot sampled
after the exposure.  For each choice of other factors the
exceptional proportion in $t_\ell$ is $o(1)$ with a uniform bound.
Integrate first against the independent uniform dyadic/residue
laws and then use the conditional domination already proved.
Each removal therefore fails with probability $o(1)$.  A union
bound over the finitely many factors and the triangle inequality
give agreement, to error tending to zero in probability, between
the original marginal cube law and the law with only its
$d$-residue restrictions.  The same reasoning applies to a test
composed with the fixed face automorphism; its Lipschitz bound is
still fixed.

In this latter law, applying $\mathcal S_e$ on the upper face
$\omega_j=1$ replaces $x_j$ by $x_j+\Delta_e$ in
\eqref{eq:alignment-joint-polynomial-exponent}, at all slots and
for all observables of this component simultaneously.
By \eqref{eq:alignment-residue-divisibility}, $d\mid\Delta_e$
eventually.  Moreover $|\Delta_e|/Z\to0$ uniformly, since its
size is bounded by a fixed power of $M+D_i^-$ whereas $R$
dominates all those powers.  Translation of the corresponding
box changes normalized counting averages by
$O(|\Delta_e|/Z+d/Z)=o(1)$.  Comparing on both sides establishes
the asserted face invariance.  Taking the finite union over
targets and own pairs is harmless.  This argument proves marginal
invariance; it makes no assertion about the other components.

\paragraph{The countable diagonal and the period.}
Enumerate the conditions consisting of a positive integer modulus,
a cube dimension at most $s+1$, bounds $C\in\N$, lower side
bounds $c=1/j$, Lipschitz bounds in $\N$, and tolerances $1/j$.
The assertion just proved holds uniformly over the uncountably
many endpoints and tests within each such condition.  Each of
the first finitely many conditions fails with probability tending
to zero.  Given events of probability at least $\epsilon>0$ along
a subsequence, choose an increasing number of initial conditions
whose union of failures has probability less than $\epsilon/2$,
and select one conditioning from the event outside that union.
By making the initial list and the accuracy increase successively,
we obtain a sequence satisfying every fixed condition with error
tending to zero.

Now fix the finite choices of ambient groups along a subsequence
and apply Proposition~\ref{prop:local-cube-law}.  Its rational
filtration, period $d_0$, and rational representatives may depend
on this sequence.  The face invariances have already been secured
for every fixed positive integer modulus, so they hold for $d_0$.
For any fixed small relative box size $\alpha>0$, they hold for
base boxes near a given macroscopic point $\beta$ and increment
boxes near zero, with the residue restrictions specified in that
Proposition.  First take the sequence limit and then let
$\alpha\downarrow0$.  The projected local cube Haar measure is
therefore invariant under the face automorphism.  Since it has
full support on its compact cube coset and the automorphism is a
homeomorphism, that coset is preserved.  This proves the last
assertion with the required order of limits.
\end{proof}

\subsection{From the word plan to positive conditional mass}

Fix one of the finitely many incoming scale tuples at pivot $i$.
Lemma~\ref{lem:finite-word-plan} gives finitely many complete
options, each recording a cumulative jump word, a resulting scale
tuple, and its slot.  Enlarge their number to one common bound
$Q_i\ge1$, and let $V_i\ge1$ bound the number of possible slots,
uniformly over incoming scales.  These constants depend only on
$n,r,s$.

For given earlier variables and an integer pivot value $x\in\N$, let
$\mathcal P_i(b;x)$ be the following event: some allowed outcome
$b'$ from the incoming scale $b$ satisfies, for every target
$B=T\cup\{i\}$, every $d\in\mathcal D_i$, color $c$, and
$D\subset\mathcal E_B$,
\begin{equation}\label{eq:alignment-pivot-property}
 S_{B,b'_Bd_B,c}(t_Tx)>2\tau
 \ \Longrightarrow\ %
 S_{B,b'_Bd_B,c}\left(t_Tx+
       \sum_{A\in D}\frac{h_A b'_A d_A}{h_B b'_B d_B}t_A\right)>\tau.
\end{equation}
For the large $w$ under consideration, every displayed offset is
an integer.  This follows either from admissibility and the finite
rational lists, or directly from the integer powers and
\eqref{eq:alignment-residue-data}.

\begin{lemma}[Conditional alignment mass]\label{lem:alignment-mass}
For every fixed $\tau>0$ and every permitted bounded-complexity
family of models, outside a set of earlier-variable/microcell/residue
conditionings of probability tending to zero, the conditional
uniform pivot probability that
\[
 \mathcal P_i(b;p+s(v))\quad\text{holds for some prescribed slot }v
\]
is at least $1/(2Q_i)$.  The same statement holds for the conditional
harmonic probability, after replacing the constant by $1/(3Q_i)$
if necessary.  The constants are independent of the threshold and
model complexity.
\end{lemma}

\begin{proof}
Suppose the uniform assertion fails on sets of conditionings of
probability bounded below along a subsequence.  By
Lemma~\ref{lem:conditional-face-invariance}, choose a sequence in
those sets with all its empirical face invariances.  Pass to a
subsequence fixing the ambient array nilmanifolds and using
Proposition~\ref{prop:local-cube-law}.  There are only finitely many
observable formulae in a conditioning.  Their common supremum
and Lipschitz bounds give, by compactness, a further subsequence
on which all of them converge uniformly.  Denote the limiting
observables by $F_{B,a,c,v}$.

For each stored factor and integer input $u$, define the reading of a
joint state $y\in K/\Lambda$ by
\begin{equation}\label{eq:alignment-intrinsic-reading}
 \mathcal R_{B,a,c;v,u}(y)
 =F_{B,a,c,v}\bigl(\operatorname{ev}_u\pi_{B,a,c,v}y\bigr).
\end{equation}
Here $\pi_{B,a,c,v}$ selects the indicated array factor and
$\operatorname{ev}_u$ evaluates it at the integer input $u$.
The coordinates of $u$ are indexed by the own pairs $(B,A)$;
when adding $u$ to a slot vector, insert zeros at all other pairs.
These readings are continuous functions of $y$ and are defined
before choosing any lifts on a point-law coset.

Consider a point-law coset of that Proposition,
\[
 h(\beta)H\sigma_r\Lambda/\Lambda,
\]
with its Haar image measure $\mu_{\beta,r}$.  The preceding Lemma
and Proposition~\ref{prop:face-lift} lift each own-pair array
shift to a polynomial shear of the joint cover coordinates
$z\in\mathfrak h$.  Such a lift may move other target components.
Adjoin its translation $v\mapsto v+\mathbf e_e$ and call the
result $L_e$.  Lemma~\ref{lem:shear-haar} puts these transformations
in a nilpotent group of step at most $s$, including the ordinary
slot translations.  Write $\Theta(z)=h(\beta)\exp(z)\sigma_r\Lambda$.
At slot $v$, color a state by the palette
\begin{equation}\label{eq:alignment-limit-palette}
 C_{B,a}(z,v)=
 \bigl(\1_{\{\mathcal R_{B,a,c;v,0}(\Theta(z))>3\tau/2\}}\bigr)_{c\in[r]}.
\end{equation}
These palettes have at most $2^r$ values, as required by the finite plan.

Apply that plan pointwise starting at any $(z,0)$.  It returns
one of the $Q_i$ options, a jumped center, and terminal palette
equalities there.  Only own letters occur in a terminal word for
$B$.  Write $m=(m_e)_e$ for their integer powers.  Their projection on
this entire component is exactly the commuting array shift
$u\mapsto u+(m_e)_e$, while their slot effect is
$v\mapsto v+m$, with the zero extension just specified.
Consequently at a center state $y$ and slot $v$ the two readings
are $\mathcal R_{B,a,c;v,0}(y)$ and
$\mathcal R_{B,a,c;v+m,m}(y)$.  This statement holds even if earlier
jumps involved letters belonging to other targets.

To check the numerical meaning, use the corresponding prelimit
observables and evaluate the same readings on an actual array state
$Y(\ell)$, with $p=p_*+M\ell$.  The center value is the
model at $t(p+s(v))$.  The translated value is the model at
\begin{align}
 t\left(p+s(v)+\sum_e m_er_e\right)
       +M\sum_e m_e\Delta_e
 &=t(p+s(v))+\sum_e m_eq_e.
       \label{eq:alignment-own-letter-reading}
\end{align}
Each occurrence uses the correct residue-specific observable at
its new slot.  For the terminal word at outcome $b'$ and future
multiplier $d$, the powers are
\[
 m_e=q_0\frac{b'_A}{b'_B}\frac{d_A}{d_B}\quad(e=(B,A),\ A\in D),
 \qquad m_e=0\quad(A\notin D).
\]
They are integers by the plan, and
$m_eq_e=(h_A b'_A d_A/(h_B b'_B d_B))t_A$.
The leading index is $b'_Bd_B$.  Therefore
\eqref{eq:alignment-own-letter-reading} gives exactly
\eqref{eq:alignment-pivot-property} with $x=p+s(v)$.
This numerical interpretation is used for terminal own-letter
words.  An arbitrary cumulative jump need not be a translation
of the actual pivot orbit.

For an option $o$ with scale outcome $b'$ and slot $v$, let
$F_o\subset K/\Lambda$ consist of the states $y$ satisfying every
required comparison.  Explicitly, for each $B,d,c,D$, put
$a=b'_Bd_B$ and take $m$ as above, and impose
\begin{equation}\label{eq:alignment-closed-comparison}
 \mathcal R_{B,a,c;v,0}(y)\le3\tau/2\quad\text{or}\quad
 \mathcal R_{B,a,c;v+m,m}(y)\ge3\tau/2.
\end{equation}
These finitely many conditions are closed by continuity of the
readings.  They depend only on the state,
the option's finite data, and the limiting formulae.  In
particular they do not depend on the choice of lifts.
For every starting cover coordinate, the pointwise palette
construction puts the $K/\Lambda$-valued coordinate of the jumped
state into $F_o$ for at least one option $o$. Equality of palettes implies
\eqref{eq:alignment-closed-comparison}, including at equality
with the threshold.

Sample $z$ from the expanding weighted boxes of
Lemma~\ref{lem:shear-haar}. For each fixed option $o$, let $Y_o$ be the
$K/\Lambda$-valued coordinate of the state after its cumulative jump.
The jump acts on the cover coordinate $z$ by a fixed shear, so the law
of $Y_o$ converges to $\mu_{\beta,r}$. Pointwise coverage gives
\[
 1\le\sum_o\Pr\{Y_o\in F_o\}.
\]
Using the closed-set direction of Portmanteau separately for the
finitely many summands yields
\begin{equation}\label{eq:alignment-haar-positive-mass}
 1\le\sum_o\mu_{\beta,r}(F_o),\qquad
 \mu_{\beta,r}\left(\bigcup_o F_o\right)\ge\frac1{Q_i}.
\end{equation}
The last implication also follows by summing the indicator
inequality $\sum_o\1_{F_o}\le Q_i\1_{\cup_oF_o}$.
This is the only quantitative loss in passing from pointwise
recurrence to the point-law measure.  Although the shears may
depend on the coset, the sets $F_o$ do not.  Therefore
\eqref{eq:alignment-haar-positive-mass} integrates over the
macroscopic-position/residue mixture without a measurable choice
of lifts.

Replace each comparison \eqref{eq:alignment-closed-comparison}
by the open comparison
\begin{equation}\label{eq:alignment-open-comparison}
 \text{center}<7\tau/4\quad\text{or}\quad
 \text{translated value}>5\tau/4,
\end{equation}
and take the same finite intersections and union.  This defines
an open set containing $\bigcup_oF_o$, so it has mixture measure
at least $1/Q_i$.  Proposition~\ref{prop:local-cube-law} gives
weak convergence of the actual uniform conditional state laws
to that mixture.  Open-set Portmanteau therefore gives lower
limit at least $1/Q_i$ for this open event.
For sufficiently late terms, uniform convergence of all
observables has error less than $\tau/4$.  If the actual center
is greater than $2\tau$, its limiting reading is greater than
$7\tau/4$, forcing the second alternative in
\eqref{eq:alignment-open-comparison}; the actual translated
reading is then greater than $\tau$.  The good boundary condition
ensures that these actual readings use exactly the selected
model formulae.  By \eqref{eq:alignment-own-letter-reading} the
open event implies success at some $p+s(v)$.

This contradicts the selected conditional probabilities being
less than $1/(2Q_i)$.  Hence the exceptional sets have probability
tending to zero.  The uniform-to-harmonic total variation estimate
following \eqref{eq:alignment-microcell} proves the last statement.
The compactness subsequence and how late the estimates hold may
depend on $\tau$ and the fixed model complexity.  The constants
$Q_i$ and the displayed lower bounds do not.
\end{proof}

\subsection{Returning to the raw variables and successive pivots}

The conditional conclusion permits a bounded list of shifts of
the pivot.  We now remove those shifts before proceeding to the
next pivot.  This step conditions only on earlier raw variables;
translation invariance is not claimed inside a fixed microcell.

Let $T_a$ denote translation by $a$.  For the raw pivot law $\mu_i$,
put $X=X_i$.  Lemma~\ref{lem:sampling}(2) gives, for $a\in W\Z$
with $|a|=o(X)$,
\begin{equation}\label{eq:alignment-raw-translation}
 \norm{T_a\mu_i-\mu_i}_{\TV}
 \le\frac{|a|}{X(\log X-W/X)}
 \le (1+o(1))\frac{|a|}{X\log X}.
\end{equation}
Here probability total variation is one half of the total-mass norm
used in Lemma~\ref{lem:sampling}.

In our application each prescribed slot satisfies
\[
 W\mid s(v),\qquad |s(v)|\le C MW^w\le CM^2=o(X_i),
\]
uniformly in all earlier raw variables.  Thus
\eqref{eq:alignment-raw-translation} gives $o(1)$ uniformly for
every such slot.  In particular it remains valid after integrating
against an arbitrary event depending on earlier variables.

\begin{proof}[Proof of Principle~\ref{pr:alignment}]
For the given $n,r,s$ take the lists and all budgets from
Lemma~\ref{lem:finite-word-plan}.  Let $I$ be the set of pivots
having targets and define
\begin{equation}\label{eq:alignment-delta}
 \delta=\prod_{i\in I}\frac1{3Q_iV_i}>0,
\end{equation}
with the empty product equal to $1$.  This definition precedes
the choice of admissible parameters, models, and $\tau$.
Now fix any such parameters, any family of models of step at most
$s$ and bounded complexity in the sense of
Definition~\ref{def:piecewise-model}, and any fixed $\tau>0$.

Suppose $E$ is an event of the earlier raw variables on a given
finite path, with incoming scale $b$ and all the earlier pivot
requirements secured.  Lemma~\ref{lem:alignment-mass}, integrated
over the pivot microcell and residue, implies
\begin{equation}\label{eq:alignment-extension-shifted}
 \frac1{3Q_i}\PP(E)-o(1)
 \le\sum_v\PP\bigl(E,\mathcal P_i(b;p+s(v))\bigr).
\end{equation}
The exceptional set of conditionings has probability $o(1)$
unconditionally, so its intersection with $E$ also costs $o(1)$.
Conditional on the earlier variables each $s(v)$ is fixed.
Equation~\eqref{eq:alignment-raw-translation} therefore changes
each summand by $o(1)$ when $p+s(v)$ is replaced by $p$.  There
are at most $V_i$ slots, and consequently
\begin{equation}\label{eq:alignment-extension}
 \PP\bigl(E,\mathcal P_i(b;p)\bigr)
 \ge\frac1{3Q_iV_i}\PP(E)-o(1).
\end{equation}
This estimate is uniform over the finite incoming-scale list and
the finite path events.  No conditional density assumption on
$E$ has been used.

On the successful extension choose the first successful scale
outcome in a fixed enumeration.  This is a measurable choice
from a finite list and partitions that event into finitely many
new path events.  They depend only on raw variables through the
current pivot.  Apply \eqref{eq:alignment-extension} to each path
at the next pivot.  Pivots without targets impose no condition
and can be skipped.  Summing the finitely many path estimates
and iterating yields total success probability at least
$\delta-o(1)$.

For completeness, consider explicitly a requirement secured at
pivot $i$ when its resulting scale was $b^{(i)}$.  Every later
update multiplies the scale by one of the finite tuples used to
construct $\mathcal D_i$.  On any ensuing path, their complete
product is some $d\in\mathcal D_i$.  Requirement
\eqref{eq:alignment-pivot-property} at pivot $i$ was required
for this very $d$, at the original earlier raw variables and
the original raw pivot.  Later steps change none of those raw
variables.  Therefore at the final scale
$b^{\rm fin}=b^{(i)}d$ its leading index is exactly
$b^{\rm fin}_B$, and every offset ratio is exactly
$h_A b^{\rm fin}_A/(h_B b^{\rm fin}_B)$.
This verifies preservation between pivots in addition to the
within-pivot preservation proved in
\eqref{eq:alignment-leading-budget}--\eqref{eq:alignment-word-budget}.

Every successful final path thus gives $b^{\rm fin}\in\mathcal B$
such that, simultaneously for every block $B$, color $c$, and
$D\subset\mathcal E_B$,
\[
 S_{B,b^{\rm fin}_B,c}(t_B)>2\tau
 \quad\Longrightarrow\quad
 S_{B,b^{\rm fin}_B,c}\left(t_B+
       \sum_{A\in D}\frac{h_A b^{\rm fin}_A}
                             {h_B b^{\rm fin}_B}t_A\right)>\tau.
\]
For blocks with no adding pairs the sole comparison is the
identity, which holds automatically.  Taking the ordinary lower
limit as $w\to\infty$ gives probability at least $\delta$.
The lists and $\delta$ depend only on $n,r,s$, whereas the rate
of convergence may depend on the fixed threshold and model
family.  These are exactly the quantifiers of
Principle~\ref{pr:alignment}.
\end{proof}
